\section{A local gap limit and distant-label decorrelation in the CUE limit}

\begin{theorem}[NN-A016: periodic gap displacement and pair decorrelation]
Let $D_j$ be the $j$th positive point of the reduced Palm sine process
at zero, with kernel $K^0$ as in NN-A003. It has a density $f_j$.
For $j\ge2$ put $\sigma_j=\sqrt{\log j}/\pi$ and extend $f_j(s)$
by zero for $s\le0$. Then
\[
 \int_{\mathbb R}\left|
 \sigma_j f_j(j+\sigma_j y)
       -\frac{e^{-y^2/2}}{\sqrt{2\pi}}\right|dy\longrightarrow0.
                                                        \tag{A016.1}
\]
In particular, for every fixed nonzero real $a$,
\[
 \mathbb E e^{ia(D_j-j)}\longrightarrow0,               \tag{A016.2}
\]
and $(D_j-j)\bmod2$ tends to Haar measure on $\mathbb R/2\mathbb Z$.
Consequently the limiting CUE label correlations in NN-A003 satisfy
\[
 \mathbb E[1-2d_2(j-D_j)]\longrightarrow0.
\]
Thus two labels far apart in the limiting indexed model become
independent fair signs. The order of limits is first $N\to\infty$
for fixed $j$ in NN-A003, then $j\to\infty$ here. No uniform
simultaneous-$j,N$ claim or arithmetic transfer is asserted.
\end{theorem}

\begin{proof}
We include the local-limit step, because an ordinary central limit
theorem alone would not imply (A016.2).

Write $k(x)=\sin(\pi x)/(\pi x)$. The sine kernel $K(x,y)=k(x-y)$
is the orthogonal projection in $L^2(\mathbb R)$ onto Fourier
frequencies $[-1/2,1/2]$. Its compression $A_s$ to $(0,s)$ is a
positive trace-class contraction with trace $s$. The determinant
generating function gives
\[
 \mathbb E z^{N(0,s)}=\det(I+(z-1)A_s)
       =\prod_\ell(1-\lambda_\ell+\lambda_\ell z).
\]
Thus this count has the distribution of a sum of independent Bernoulli
variables with parameters the eigenvalues $\lambda_\ell\in[0,1]$.
The product converges because $\sum\lambda_\ell=s<\infty$;
its probability generating function identifies the distribution.
Direct integration gives
\[
 \operatorname{Var}N(0,s)
 =s-2\int_0^s(s-u)k(u)^2du
 =\frac{\log s}{\pi^2}+O(1).                          \tag{A016.3}
\]
Indeed $\int_0^\infty k(u)^2du=1/2$; the remaining tail
$2s\int_s^\infty k(u)^2du$ is bounded, while
$2\int_0^s uk(u)^2du=(\log s)/\pi^2+O(1)$ by
$\sin^2(\pi u)=(1-\cos(2\pi u))/2$.

Conditioning on a fixed point removes its kernel evaluation line
from the projection. Conditioning on two distinct points $0,s$
removes their two-dimensional span. This follows either directly
from the Schur complement of the determinant correlation functions
or from the Gram determinant formula for those evaluation vectors.
Let $B_s$ be the restriction to $(0,s)$ of this two-point reduced
Palm projection. Then $A_s-B_s\ge0$ has trace at most two and rank
at most two. Hence
\[
 \operatorname{Tr}B_s=s+O(1),\qquad
 \operatorname{Tr}(B_s-B_s^2)
       =\frac{\log s}{\pi^2}+O(1).                    \tag{A016.4}
\]
For clarity, if $R=A_s-B_s$, cyclicity of trace gives the variance
difference $\operatorname{Tr}R-\operatorname{Tr}(R(A_s+B_s))$.
Since $0\le A_s+B_s\le2I$, its absolute value is at most
$\operatorname{Tr}R\le2$. The same assertions, with rank one,
hold for the Palm process at zero alone. All these bounds are
uniform as $s\to\infty$.

Here is the needed local central limit lemma. If $X$ is any finite
or countable sum of independent Bernoulli variables, has finite
mean $\mu$, and variance $v\to\infty$ in a sequence of such sums,
then
\[
 \sup_{k\in\mathbb Z}\left|
 \sqrt v\,\mathbb P(X=k)
 -\frac1{\sqrt{2\pi}}e^{-(k-\mu)^2/(2v)}
 \right|\longrightarrow0.                            \tag{A016.5}
\]
This statement is uniform over the Bernoulli parameters once $v$
is sufficiently large. To prove it, the centered characteristic
function $\psi$ satisfies, for $|\theta|\le\pi$,
\[
 |\psi(\theta)|\le
 \exp[-2v\sin^2(\theta/2)]
 \le\exp[-2v\theta^2/\pi^2].
\]
For fixed $u$, Taylor expansion of the Bernoulli logarithms gives
$\log\psi(u/\sqrt v)=-u^2/2+O(|u|^3/\sqrt v)$.
The remainder is uniform because the third derivative of each
centered Bernoulli logarithm, for small argument, is bounded by a
constant times $p(1-p)$. Fourier inversion, scaling by $\sqrt v$,
and dominated convergence on the interval
$[-\pi\sqrt v,\pi\sqrt v]$ prove (A016.5): the uniform error
in $k$ is at most the $L^1$ difference between the scaled
characteristic function, extended by zero, and $e^{-u^2/2}$.
The same argument applies to the countable product by finite
truncation. This proves the lemma without a local-limit assumption.

The Palm Campbell formula gives a density for the $j$th positive
point:
\[
 f_j(s)=K^0(s,s)\,
 \mathbb P^{0,s}\{N(0,s)=j-1\},\qquad s>0,            \tag{A016.6}
\]
where $\mathbb P^{0,s}$ is the two-point reduced Palm law.
Indeed summing a test function over points $s$ with exactly $j-1$
points between $0$ and $s$ selects precisely the $j$th point;
conditioning that summand at $s$ gives the right side.
The count in (A016.6) has the Bernoulli eigenvalue description for
$B_s$. Also $K^0(s,s)=1-k(s)^2\to1$.
For $s=j+y\sigma_j$ with $y$ in any fixed compact set, (A016.4)
and (A016.5) therefore give, uniformly in $y$,
\[
 \sigma_j f_j(j+y\sigma_j)
       \longrightarrow\frac{e^{-y^2/2}}{\sqrt{2\pi}}. \tag{A016.7}
\]
Only an $O(1)$ error in the Palm mean is needed, since
$\sigma_j\to\infty$.

To control tails, apply Chebyshev to the one-point Palm count,
whose mean is $s+O(1)$ and variance is (A016.3) up to $O(1)$.
For each fixed large $M$, the identities
$D_j>s\iff N(0,s)\le j-1$ and their lower-tail counterparts give
\[
 \limsup_{j\to\infty}
 \mathbb P(|D_j-j|>M\sigma_j)\le C/M^2.
\]
There are infinitely many positive points almost surely: the same
variance estimate and Chebyshev at dyadic $s$ give this by
Borel--Cantelli. Thus (A016.6) has total integral one.
Compact-uniform density convergence (A016.7), the displayed tail
bound, and Gaussian tail decay prove (A016.1).

Finally write $g_j(y)=\sigma_j f_j(j+\sigma_j y)$ and
$\varphi(y)=e^{-y^2/2}/\sqrt{2\pi}$. For $a\ne0$,
\[
 \left|\mathbb E e^{ia(D_j-j)}-e^{-a^2\sigma_j^2/2}\right|
 \le\|g_j-\varphi\|_{L^1}\longrightarrow0.
\]
This proves (A016.2), and taking $a=\pi k$ proves all nonconstant
Fourier coefficients of the residue-class law vanish. The continuous
periodic function $1-2d_2(x)$ has Haar mean zero, so its expectation
tends to zero. In NN-A003 both signs have zero mean and the joint
law is unchanged by flipping both, because the uniform phase can
be shifted by one. Their four joint probabilities are therefore
$(1\pm C_j)/4$, where $C_j$ is that correlation, and tend to $1/4$.
\end{proof}

The geometry in this result is again the sine projection, finite-rank
Palm modifications, and positive trace-class contractions on intervals.
The full Fredholm generating function supplies local count information
that a pair-correlation statement does not contain. The proof illustrates
the specific strengthening needed to pass from a Gaussian scaling law
to a fixed-frequency periodic statistic. It asserts neither the
arithmetic counterpart nor mixing of every finite label block.
